\chapter{Normal sources and supported operations}
\label{mht:ch:normal-sources}

A closed field coupled to a boundary algebra must do more than
produce a scalar cyclic cochain. Its insertion has a normal
support, a differential and collision operations. The action
determines a source in the closed equation as well as an action
on the open fields. We first construct a classical model in
which these two effects are transposes of one another.

Normal jets then give explicit supported cyclic operations.
Real normal currents, holomorphic principal parts and smooth
source tests require different complexes. We compute their
primitives and the obstruction to extending an operation across
a collision. These constructions retain the data needed for
a comparison with independently defined closed fields.
The final section studies coefficient completion. It proves
reconstruction in one normal direction and identifies the
limits and duals that fail to preserve support.

\section{A coupled Hamiltonian bulk and defect}

Let
\[
 M=\mathbb R_r\times\mathbb R_s\times\mathbb C^2_{q,p},
 \qquad Y=\{r=0\},\qquad i:Y\hookrightarrow M,
 \qquad \Omega=dq\wedge dp.
\]
Orient the real normal line by increasing $r$. It is the whole
line, so $Y$ is an internal defect.

Let $\mathcal A_M$ denote
smooth mixed forms, with the completed projective smooth
topology and differential
\[
 D_M=d_r+d_s+\bar\partial.
\]
Define $\mathcal A_Y$ similarly, with $D_Y=d_s+\bar\partial$.
An integral always has compact support in its integration
variables. An unintegrated smooth factor may act as a multiplier.

Complete coefficientwise in $\epsilon$, of degree zero. The
holomorphic Moyal product is
\begin{equation}\label{mht:eq:geometric-moyal}
 a\star b=m\exp\!\left(\frac\epsilon2\mathsf P\right)(a\otimes b),
 \qquad
 \mathsf P=\partial_q\otimes\partial_p
                    -\partial_p\otimes\partial_q.
\end{equation}
Exterior and matrix multiplication are included in $m$.
Every parameter coefficient is a finite-order differential
operator.

For scalar mixed forms the graded commutator is
divisible by $\epsilon$. Hence
\[
 [f,g]_\epsilon=\epsilon^{-1}[f,g]_\star,
 \qquad
 \rho_f(a)=\epsilon^{-1}[i^*fI_N,a]_\star
\]
are defined before parameter inversion. The second acts on
matrix-valued defect forms. Associativity gives
\[
 [\rho_f,\rho_g]=\rho_{[f,g]_\epsilon},\qquad
 \rho_f(a\star b)=(\rho_fa)\star b
                    +(-1)^{|f||a|}a\star\rho_fb.
\]

The trace $\tau_Y(a)=\int_Y\Omega\wedge\operatorname{Tr}a$
is graded cyclic. Integration by parts kills every positive
two-input Moyal correction, so $\tau_Y(a\star b)=\tau_Y(ab)$.
The action $\rho$ preserves the resulting pairing. Here $q,p$
are physical coordinates. Proposition~\ref{mht:prop:physical-spencer}
compares their jets with horizontal independent coefficients.
Unconstrained coefficient fields have a different carrier.

Let $\mathcal A_M^{\mathrm{nor}}$ admit distributions of locally
finite order in $r$, with smooth distribution-valued tangential
dependence. It is a module over smooth mixed forms and is
stable under $D_M$. In particular it contains every finite
derivative of $\delta(r)$.

We use the fields
\[
 \alpha\in\mathcal A_M[1][[\epsilon]],\qquad
 \beta\in\mathcal A_M^{\mathrm{nor}}[2][[\epsilon]],\qquad
 A\in(\mathcal A_Y\otimes\operatorname{Mat}_N)[1][[\epsilon]].
\]
The bulk cotangent pairing and the defect trace pairing have
degree $-1$. On the distributional domain the pairing is
weak: we use its Hamiltonian bracket on functionals whose
variational derivatives are represented in the paired field
spaces. The functionals below have this property.

At ghost number zero, $\alpha,A,\beta$ have
mixed form degrees $1,1,2$. Define
\begin{align}
 S_c&=\int_M\Omega\wedge\beta\star
             \left(D_M\alpha+\tfrac12[\alpha,\alpha]_\epsilon\right),
       \label{mht:eq:geometric-bulk-action}\\
 S_o&=\tau_Y\left(\tfrac12 A\star D_YA
                             +\tfrac13 A\star A\star A\right),
       \label{mht:eq:geometric-open-action}\\
 S_{co}&=\tfrac12\tau_Y(A\star\rho_\alpha A).
       \label{mht:eq:geometric-coupling}
\end{align}
For all BV components, use suspension of degree $-1$ with
$b_1(sx)=-sDx$ and $b_2(sx,sy)=(-1)^{|x|}s[x,y]$.
The cyclic pairing fixes the signs of the corresponding
quadratic and cubic Hamiltonians.

\begin{theorem}\label{mht:thm:geometric-coupled-classical}
The functional $S_c+S_o+S_{co}$ satisfies the classical
master equation. Its Hamiltonian vector field preserves the
smooth and normal-distribution domains coefficientwise in
$\epsilon$. It uses no product of two singular normal
coefficients.
\end{theorem}
\begin{proof}
Let $\mathfrak g=(\mathcal A_M[[\epsilon]],D_M,[-,-]_\epsilon)$,
and let $\mathfrak a$ be the matrix defect dg Lie algebra with
its graded star commutator. Its cyclic pairing has degree $-3$.
Adjoin to $\mathfrak g\oplus i_*\mathfrak a$ the admissible
paired dual of $\mathfrak g$. Integration identifies this
summand with $\mathcal A_M^{\mathrm{nor}}[1][[\epsilon]]$
before the common suspension.

The dual summand is abelian. Let $\mathfrak g$ act on it
coadjointly and on $\mathfrak a$ by $\rho$. Its bracket
with $\mathfrak a$ is zero. The bracket of two open elements
has its native component and a dual component $\mu$, determined
by cyclicity:
\[
 \langle\mu(a,b),x\rangle=\langle a,[b,x]\rangle.
\]
For degree-zero inputs this reads
$\mu(a,b)(x)=\langle\rho_xa,b\rangle$.
The transpose of restriction is normal pushforward.

Jacobi on two bulk inputs and one open input is the
representation identity for $\rho$. On one bulk and two
open inputs it is the derivation identity and equivariance
of $\mu$. For three open inputs the native component is
Jacobi in $\mathfrak a$. Pairing the dual component with
$x$ gives the cyclic sum of $\langle\rho_x[a,b],c\rangle$.
The derivation identity and cyclicity cancel its terms with
their Koszul signs. The remaining triples give coadjoint
identities or zero.

The mixed differentials preserve these brackets by integration
by parts. We have therefore constructed a cyclic dg Lie
algebra. Its suspended Hamiltonian is exactly the sum of
\eqref{mht:eq:geometric-bulk-action}--\eqref{mht:eq:geometric-coupling},
and the dg Lie identities give the classical master equation.
All singular coefficients lie in the abelian dual summand.
Differentiation and multiplication by smooth fields preserve
that summand, proving the domain assertion.
\end{proof}

The source follows directly by varying the coupling. Since
$A$ is odd, $\operatorname{Tr}(A^2)=0$ at $\epsilon=0$.
Graded exchange removes the even powers of $\mathsf P$, giving
\begin{equation}\label{mht:eq:geometric-source-moment}
 \mathcal M_\epsilon(A)=\epsilon^{-1}\operatorname{Tr}(A\star A)
 =\sum_{k\geq0}\frac{\epsilon^{2k}}{2^{2k+1}(2k+1)!}
       \operatorname{Tr}m\mathsf P^{2k+1}(A\otimes A).
\end{equation}
Its leading term is $\operatorname{Tr}(\partial_qA\wedge\partial_pA)$.

Cyclic rotation rewrites the coupling as
$\int_Y\Omega\wedge i^*\alpha\star\mathcal M_\epsilon(A)$.
The ghost-zero equations are
\begin{align}
 D_M\alpha+\tfrac12[\alpha,\alpha]_\epsilon&=0,\notag\\
 D_YA+A\star A+\rho_\alpha A&=0,\notag\\
 D_M\beta+[\alpha,\beta]_\epsilon
       &=\delta(r)dr\,\mathcal M_\epsilon(A).
       \label{mht:eq:geometric-source-equation}
\end{align}
The last sign uses $dr$ before the orientation of $Y$.
Equivalently, variation of $S_{co}$ in $\alpha$ has the
positive pairing with $\mathcal M_\epsilon(A)$.

The source is covariantly closed on an open solution.
Indeed $D_\alpha=D_Y+\rho_\alpha$ is an odd derivation and
$D_\alpha A=-A\star A$, so
\[
 D_\alpha(A\star A)
   =(-A\star A)\star A-A\star(-A\star A)=0.
\]
Trace intertwines the scalar Hamiltonian action. This is
the source condition in the bulk Bianchi identity.

For two coincident blocks, the source also contains
\[
 \epsilon^{-1}\operatorname{Tr}
       (A_{12}\star A_{21}+A_{21}\star A_{12}).
\]
The cubic action retains the corresponding mixed cyclic
words. A block-diagonal field is a subtheory of this
rank-added defect, not its full field space.

\section{Inner actions and normal transport}

An inner action can be removed from the open differential
while leaving a closed cubic interaction. We first explain
the algebraic change of variables. Let $\mathfrak a$ be an
ordinary quadratic Lie algebra, and suppose
$\rho_x=\operatorname{ad}_{j(x)}$ for a Lie map
$j:\mathfrak g\to\mathfrak a$.
Write $\kappa(x,y)=\langle jx,jy\rangle$ and let
$j^\dagger$ be the adjoint into the paired dual.

\begin{proposition}\label{mht:prop:inner-double-extension}
In the double extension, the change
\[
 a'=a+jx,\qquad
 \xi'=\xi-j^\dagger a-\tfrac12j^\dagger jx
\]
preserves the pairing and decouples the open Lie algebra.
The remaining bulk-dual bracket has the cocycle
$c(x,y)=\kappa([x,y],-)/2$.
Its cyclic cubic Hamiltonian is $\kappa(x,[x,x])/12$.
\end{proposition}
\begin{proof}
Substitution in $\langle a,a\rangle+2\langle x,\xi\rangle$
cancels both cross terms and the term $\kappa(x,x)$.
The new open vector has old coordinates $(0,a,j^\dagger a)$.
The new bulk vector has coordinates
$(x,-jx,-j^\dagger jx/2)$.

Insert these vectors in the double-extension bracket.
The mixed bracket vanishes. The open bracket is the
original bracket of $\mathfrak a$, and the bracket of two
new bulk vectors is $([x,y],0,c(x,y))$.
Invariance of $\kappa$ proves the cocycle identity.
Pairing with the third bulk vector and multiplying by
$1/6$ gives the cubic Hamiltonian. The same calculation
with the cyclic Koszul signs applies after suspension.
\end{proof}

In the geometric theory $j(\alpha)=\epsilon^{-1}i^*\alpha I_N$
requires inversion of $\epsilon$. The canonical change is
\[
 A'=A+\epsilon^{-1}i^*\alpha I_N,\qquad
 \beta'=\beta-\delta(r)dr
       \left(\epsilon^{-1}\operatorname{Tr}A
                         +\frac{N}{2\epsilon^2}i^*\alpha\right).
\]
It gives
\begin{equation}\label{mht:eq:surviving-cyclic-cubic}
 S=S_c(\alpha,\beta')+S_o(A')
       +\frac{N}{6\epsilon^3}\int_Y\Omega\wedge(i^*\alpha)^{\star3}.
\end{equation}

The coefficient can be checked without suppressing the
kinetic term. Put $B=i^*\alpha I_N/\epsilon$ and
$F_B=D_YB+B\star B$. Modulo an integrated total derivative,
\[
 \operatorname{CS}(A+B)
 =\operatorname{CS}(A)+\operatorname{CS}(B)
       +\tau_Y(B\star A\star A)+\tau_Y(A\star F_B).
\]

The cotangent shifts contribute
$\tau_Y(A\star F_B)+\tau_Y(B\star F_B)/2$ when the old
$\beta$ is expressed in terms of $\beta'$. Consequently
\[
 S=S_c(\alpha,\beta')+S_o(A')
          -S_o(B)+\tfrac12\tau_Y(B\star F_B).
\]
The last two terms equal $\tau_Y(B^{\star3})/6$.
Taking the matrix trace gives
\eqref{mht:eq:surviving-cyclic-cubic}.

The residual cubic is closed and has zero self-bracket in
the cotangent theory because it does not depend on $\beta'$.
The transformation is not defined integrally at $\epsilon=0$.
Thus innerness after localization does not remove the
integral problem or its closed cubic vertex.

To move the defect, write a smooth background on a collar as
$\alpha=dr\,\alpha_r+\alpha_Y$. Complete in a field-amplitude
parameter as well as $\epsilon$. The ordered series
\begin{equation}\label{mht:eq:normal-hamiltonian-transport}
 U(r_1,r_0)=1+\sum_{n\geq1}(-1)^n
 \int_{r_0<t_1<\cdots<t_n<r_1}
 \rho_{\alpha_r(t_n)}\cdots\rho_{\alpha_r(t_1)}
                                      \,dt_1\cdots dt_n
\end{equation}
is finite at each amplitude degree. Set
\[
 D_{Y,\alpha}(r)=D_Y+\rho_{\alpha_Y(r)},\qquad
 F_{rY}=\partial_r\alpha_Y-D_Y\alpha_r
                         +[\alpha_r,\alpha_Y]_\epsilon.
\]

\begin{proposition}\label{mht:prop:normal-transport-curvature}
The operator $U$ preserves the star product and integrated
trace and composes under subdivision of the collar.
Its differential defect is
\begin{align}
 &D_{Y,\alpha}(r_1)U(r_1,r_0)-U(r_1,r_0)D_{Y,\alpha}(r_0)
       \notag\\
 &\qquad=\int_{r_0}^{r_1}
       U(r_1,t)\rho_{F_{rY}(t)}U(t,r_0)\,dt.
       \label{mht:eq:normal-transport-defect}
\end{align}
In particular a flat background gives a cyclic dg algebra map.
\end{proposition}
\begin{proof}
Differentiating the ordered simplex gives
$\partial_{r_1}U=-\rho_{\alpha_r(r_1)}U$, with initial value
the identity. The multiplicative defect satisfies the same
homogeneous equation with zero initial value, since $\rho$
is a derivation. Formal recursion proves it zero.
Trace invariance follows from $\tau_Y\rho=0$.
Subdivision of each simplex proves composition.

Call the left side of \eqref{mht:eq:normal-transport-defect}
$C(r_1)$. Differentiating gives
$\partial_{r_1}C+\rho_{\alpha_r}C=\rho_{F_{rY}}U$ and
$C(r_0)=0$. Variation of constants proves the formula.
Every step is coefficientwise finite in amplitude and
$\epsilon$.
\end{proof}

\section{Real normal jets and cyclic caps}

Normal differentiation of an insertion must act on every
input. Its algebraic model is dual to multiplication of
normal test jets. For $M\geq1$ put
\[
 J_M=\mathbb C[r,dr]/(r^{M+1},r^Mdr),\qquad
 |r|=0,\quad |dr|=1,\quad d r=dr.
\]
The second relation is the derivative of the first.
Thus $J_M$ is a finite-dimensional commutative dg algebra.

Let $C_M=J_M^\vee$ have basis $t_0,\ldots,t_M$ in degree
zero and $e_0,\ldots,e_{M-1}$ in degree $-1$.
They extract the coefficients of $r^n$ and $r^n dr$.
For a homogeneous functional use $dc=-(-1)^{|c|}c d$.

\begin{proposition}\label{mht:prop:real-normal-coalgebra}
The dual is a counital dg coalgebra with
\[
 dt_n=0,\qquad de_n=(n+1)t_{n+1},
\]
\[
 \Delta t_n=\sum_{a+b=n}t_a\otimes t_b,\qquad
 \Delta e_n=\sum_{a+b=n}(e_a\otimes t_b+t_a\otimes e_b).
\]
The counit extracts $t_0$. If $i(1)=t_0$ and $p$ is the
counit, then
\[
 Ht_n=e_{n-1}/n\quad(n>0),\qquad Ht_0=He_n=0
\]
satisfies $dH+Hd=1-ip$. The inclusions $C_M\to C_{M+1}$
preserve the differential and coproduct.
\end{proposition}
\begin{proof}
Evaluate on $d(r^k)=kr^{k-1}dr$ to obtain the differential.
Dualizing monomial multiplication gives the coproduct.
Associativity and the unit give coassociativity and the
counit identities. For $i+j=n+1$, the two contributions
to the coefficient of $t_i\otimes t_j$ in the coderivation
identity have coefficients $i$ and $j$. Their sum is $n+1$.

The formula for $H$ verifies the contraction on each basis
vector. The quotient $J_{M+1}\to J_M$ is a dg algebra map,
which proves compatibility between stages.
\end{proof}

The union $C=\varinjlim C_M$ retains every finite normal
derivative. Its iterated coproduct distributes a divided
derivative among all arguments:
\[
 \Delta^{(k-1)}t_n
   =\sum_{n_1+\cdots+n_k=n}t_{n_1}\otimes\cdots\otimes t_{n_k}.
\]
For $e_n$ the same sum has one $e$ in every possible position.
Normal differentiation raises the stage. It does not preserve
each fixed truncation.

These functionals have supported representatives. Define
\[
 E_n=\frac{(-1)^n}{n!}\delta^{(n)}(r),\qquad T_n=E_n\,dr.
\]
Then $\int T_n f=f^{(n)}(0)/n!$ and
$\int E_n g\,dr=g^{(n)}(0)/n!$.
Before shifting, $dE_n=-(n+1)T_{n+1}$.
The assignments $t_n\mapsto T_n[1]$ and $e_n\mapsto E_n[1]$
therefore identify $C$ with the finite-order supported
currents in the shifted de Rham complex.
Every evaluation is continuous in a finite smooth seminorm.

A smooth Thom form has the same normal cohomology. Choose
$\chi\in C_c^\infty(\mathbb R)$ with integral one and set
$F_\chi(r)=\int_{-\infty}^r\chi(s)\,ds$.
Before shifting, the homotopy on a compactly supported one-form is
\[
 h(g\,dr)(r)=\int_{-\infty}^r g(s)\,ds
                       -F_\chi(r)\int_{\mathbb R}g(s)\,ds.
\]
It vanishes outside a compact set. Direct differentiation
gives $dh+hd=1-ip$, with $i(1)=\chi\,dr$ and
$p(g\,dr)=\int g$. The primitive of the difference between
two unit-integral choices is compactly supported and gives
a homotopy between their inclusions.

Let $A$ now be a unital dg algebra. A test ideal can instead
use its unitalization and relative normalized chains.
In $\operatorname{CH}(A)$, an $n$-chain of internal degree
$m$ has total degree $m-n$. Set
\[
 D=d_{\mathrm{int}}+(-1)^m b,\qquad
 B_{\mathrm{tot}}=(-1)^m B.
\]
The faces defining $b$ include the graded cyclic last face.
The operator $B$ is cyclic rotation with unit insertion.

For the graded dual cochains use
\[
 (D^\vee\phi)(c)=-(-1)^{|\phi|}\phi(Dc),\qquad
 (B^\vee\phi)(c)=-(-1)^{|\phi|}\phi(B_{\mathrm{tot}}c).
\]
Their degrees are $1,-1$ and they form a mixed complex.
Completion over arities is a product.

For $A_M=J_M\otimes A$, define
\begin{equation}\label{mht:eq:normal-cap}
 \begin{split}
 \operatorname{Cap}(c,\phi)(j_0a_0,\ldots,j_na_n)
  &={}(-1)^\sigma
      \langle c,j_0\cdots j_n\rangle\phi(a_0,\ldots,a_n),\\
 \sigma&=\sum_{i<j}|a_i||j_j|+|\phi|\sum_i|j_i|.
 \end{split}
\end{equation}
Here $|\phi|$ is the total cochain degree, not its arity.

\begin{theorem}\label{mht:thm:supported-cyclic-cap}
The cap is a map of mixed complexes
\[
 C_M\otimes\operatorname{CH}(A)^\vee
       \longrightarrow\operatorname{CH}(A_M)^\vee.
\]
It preserves the cyclic and unit identities and normal-jet
refinement. It extends coefficientwise to $D^\vee+uB^\vee$,
with $|u|=2$. Continuous tangential cochains give continuous
finite-order normal operations supported at $r=0$.
\end{theorem}
\begin{proof}
Move all normal factors to the left and multiply them.
This gives a chain map
\[
 T:\operatorname{CH}(J_M\otimes A)\longrightarrow
                                      J_M\otimes\operatorname{CH}(A),
\]
with sign $\sum_{i<j}|a_i||j_j|$. The internal differential
commutes with this operation by the Leibniz rule.
Each Hochschild face has the same normal product.
The last face uses graded commutativity to move its normal
factor around that product.

The source totalization has sign $(-1)^{m_J+m_A}$ on $b$.
The tensor differential on the target has the same sign,
$(-1)^{m_J}(-1)^{m_A}$. Graded commutativity also preserves
cyclic rotation, and inserting a unit multiplies the normal
product by one. Hence $T$ commutes with $B_{\mathrm{tot}}$.
A degenerate source tensor has a unit tangential bar entry,
so the map descends to normalized chains.

Graded transposition contributes the second term of $\sigma$.
It proves the mixed identities and their descendant extension.
The formulas are unchanged under refinement. Each normal
coefficient uses finitely many smooth derivatives, proving
the support and continuity assertions.
\end{proof}

For zero-form normal inputs, the cap of $t_n$ is the sum
of tangential evaluations on all Taylor coefficients whose
orders add to $n$. If $\phi$ is closed, the cap of $e_n$
is a primitive of $(n+1)\operatorname{Cap}(t_{n+1},\phi)$.
The sign is visible for $A=\Lambda(\theta)$, $|\theta|=1$,
and the zero-arity functional $\phi(\theta)=1$.
Then $|\phi|=-1$ and
$\operatorname{Cap}(e_0,\phi)(dr\,\theta)=-1$.
Its dual differential evaluates to $+1$ on $r\theta$.

The bare scalar transposes in Chapter~\ref{mht:ch:integral-cyclic}
use a different sign convention. For an ungraded algebra,
the total-complex comparison is
\begin{equation}\label{mht:eq:scalar-dual-convention}
 u^j\phi_n\longmapsto
             (-1)^{n(n+1)/2+j}u^j\phi_n.
\end{equation}
Increasing $n$ contributes $(-1)^{n+1}$.
Decreasing $n$ and increasing $j$ contributes the same sign.
Thus \eqref{mht:eq:scalar-dual-convention} intertwines both
total differentials.

It is semilinear for $u\mapsto-u$.
It does not identify the two conventions over the identity
map of the descendant coefficient ring.

If $A$ has a closed graded trace $\tau$ and a degree-zero
derivation $X$ with $[d_A,X]=0$ and $\tau X=0$, then
$\phi_X(a_0,a_1)=\tau(a_0Xa_1)$ is a closed cyclic cochain.
The derivation identity cancels its Hochschild boundary.
Integration by parts proves cyclicity, and $X1=0$ gives
the normalized Connes identity. The cap constructs all its
finite normal descendants. A map from a physical closed
field to $\phi_X$ remains separate data.

\section{Holomorphic normal residues}

A complex normal coordinate has a Dolbeault differential.
Orient $\mathbb C_z$ by $dx\wedge dy$, so that
$d\bar z\wedge dz=2i\,dx\wedge dy$.
Define the higher polar distributions by
$z^{-n-1}=(-1)^n\partial_z^n(z^{-1})/n!$ and put
\[
 \mathcal T_n=\frac1{2\pi i}\bar\partial(z^{-n-1}).
\]
The Cauchy formula and differentiation give
\begin{equation}\label{mht:eq:holomorphic-normal-actions}
 z\mathcal T_0=0,\qquad z\mathcal T_n=\mathcal T_{n-1}\ (n>0),
 \qquad \partial_z\mathcal T_n=-(n+1)\mathcal T_{n+1}.
\end{equation}
For a test holomorphic near zero, with a compact cutoff,
$\int\mathcal T_n\wedge dz\,f=[z^n]f$.
These are complex normal residues, not the real Thom
currents of the preceding section.

For $a>0$ a smooth representative is
\[
 c_a(z)=\frac{1-e^{-|z|^2/a}}z,\qquad
 \mathcal T_{0,a}=\frac{e^{-|z|^2/a}}{2\pi i a}\,d\bar z.
\]
The numerator vanishes quadratically at zero, so $c_a$
extends smoothly. Direct differentiation gives
$(2\pi i)^{-1}\bar\partial c_a=\mathcal T_{0,a}$.
Polar integration shows that $\mathcal T_{0,a}\wedge dz$
has mass one. This fixes a one-source normalization.

The algebraic counterpart is the principal-parts module
\[
 \mathcal P_z=\mathbb C[z,z^{-1}]/\mathbb C[z]
       =\bigoplus_{n\geq0}\mathbb C[z^{-n-1}].
\]
Its residue pairing with $\mathbb C[[z]]\,dz$ is continuous
against the formal topology and separates coefficients.
Several independent normal coordinates use tensor products
of these modules. A product of distributions with the same
normal support is not supplied by this tensor construction.

\section{Supported operations and their normal contraction}

We now keep the entire observable target while resolving
the normal source differential. Let $(F,Q)$ be a Hausdorff
locally convex cochain complex with continuous differential.
No lower degree bound is assumed. Formal coefficient rings
may be completed coefficientwise.

For $n$ occupied source positions in a complex affine chart
$V$, define
\[
 \mathcal K(V)=C_c^\infty(V)\otimes\Lambda(e_1,\ldots,e_n),
 \qquad |e_i|=-1,\qquad
 d_\partial=-\sum_i\iota_{e_i}\partial_{\bar z_i}.
\]
The exterior contraction uses the displayed order of the
markers. An operation is supported on a closed set $Z$
when it kills every test supported away from $Z$.
Let $\mathcal O_Z(F)$ be the continuous supported Hom
complex, with uniform convergence on bounded source sets
in every target seminorm. Its differential is
\begin{equation}\label{mht:eq:supported-Hom-differential}
 \delta A=QA-(-1)^q A d_\partial,
 \qquad |A|=q.
\end{equation}

Choose adapted affine coordinates
$V=D\times W\subset\mathbb C^r_u\times\mathbb C^s_w$,
where $r+s=n$ and $Z=\{u=0\}$.
For a partial diagonal, $w_b$ is one representative of
each block and the other coordinates are $u_i=z_i-w_b$.
These are source collision normals. A physical normal
coordinate such as $r$ on $M$ remains inside the target $F$.

For compact $K\Subset W$, use the usual Fr\'echet space
$C_K^\infty(W)$. The normal jet spaces are
\begin{align*}
 \mathcal J_r(W)
   &=\underset{K\Subset W}{\operatorname{indlim}}
         \prod_{\mathbf a,\mathbf b\in\mathbb N^r}C_K^\infty(W),\\
 \mathcal J_r^{\mathrm{hol}}(W)
   &=\underset{K\Subset W}{\operatorname{indlim}}
         \prod_{\mathbf a\in\mathbb N^r}C_K^\infty(W).
\end{align*}
The products have their product topologies and the compact
exhaustion gives a strict LF topology. All coefficients
at one stage have a common compact support. Taylor expansion is
\[
 j\kappa=\sum_{\mathbf a,\mathbf b}
 \frac{\partial_u^{\mathbf a}\partial_{\bar u}^{\mathbf b}
                    \kappa(0,w)}{\mathbf a!\mathbf b!}
                     u^{\mathbf a}\bar u^{\mathbf b}.
\]
A normal-flat test has every displayed coefficient zero.

\begin{lemma}\label{mht:lem:bounded-normal-jets}
The Taylor map $j:C_c^\infty(V)\to\mathcal J_r(W)$ is
continuous and onto. Every bounded jet set has a bounded
set of smooth lifts. For $r>0$ there is no continuous
linear right inverse at a fixed nonzero support stage.
\end{lemma}
\begin{proof}
Derivative restriction proves continuity. A bounded jet
family lies at one support stage. Write $P_l(f;u,w)$
for the sum of its monomials of total normal degree $l$.
Choose a normal cutoff $\chi_0$ equal to one near zero.
For $m<l$, the smooth seminorm through order $m$ of
$\chi_0(u/\varepsilon)P_l$ is bounded by
$C_{l,m}\varepsilon^{l-m}$ times finitely many coefficient
seminorms. Tangential derivatives obey the same estimate.

Choose $\varepsilon_l\downarrow0$ uniformly on the bounded
family so that this bound is at most $2^{-l}$ for
$m\leq l/2$, apart from finitely many initial terms.
The series $\sum_l\chi_0(u/\varepsilon_l)P_l$ converges
uniformly in every smooth seminorm. It has common compact
support and the prescribed Taylor coefficients.

A continuous linear right inverse would bound its output
$C^0$ norm by finitely many seminorms of finitely many
input coefficients. A jet supported in a coefficient
outside this finite list would then lift to a function
with zero $C^0$ norm. This contradicts its nonzero jet.
\end{proof}

\begin{lemma}\label{mht:lem:supported-jet-factorization}
Every continuous operation supported on $Z$ factors uniquely
through normal Taylor jets, componentwise in the markers.
The resulting identification of Hom spaces preserves the
bounded-set topologies.
\end{lemma}
\begin{proof}
Let $A$ be supported on $Z$ and $\kappa$ normal-flat.
Then $A\kappa=A(\chi_0(u/\varepsilon)\kappa)$.
For each target seminorm, continuity bounds this by
finitely many test derivatives on a common compact set.
Normal flatness beats every inverse power of $\varepsilon$
from the cutoff, so $A\kappa=0$.

The Taylor map at a fixed compact stage is a surjection
of Fr\'echet spaces. The open mapping theorem makes the
factor continuous. Passing through the compact exhaustion
gives the asserted factorization. Conversely, a continuous
map on jets composed with $j$ is supported on $Z$.

Precomposition with $j$ is continuous on bounded-set Hom
spaces. For the inverse, lift each bounded jet set by
Lemma~\ref{mht:lem:bounded-normal-jets}. Taking the supremum
over these bounded lifts gives exactly the required bound
in each target seminorm.
\end{proof}

The quotient is on smooth tests. It does not restrict
a physical distribution to its own singular support.
In particular no uniform normal-order bound across all
target seminorms is required.

Change the odd basis with the affine coordinates and put
the normal markers first. The jet differential becomes
\[
 d=-\sum_{j=1}^r\iota_j\partial_{\bar u_j}
       -\sum_{j=1}^s\iota_{r+j}\partial_{\bar w_j}.
\]
For a partial diagonal, the tangent marker of a block is
the sum of its original markers. Let $\pi_j$ set
$\bar u_j=0$, and define formal integration by
\[
 S_j(u^{\mathbf a}\bar u^{\mathbf b})
  =\frac{u^{\mathbf a}\bar u^{\mathbf b+\mathbf e_j}}{b_j+1}.
\]
Thus $\partial_{\bar u_j}S_j=1$ and
$S_j\partial_{\bar u_j}=1-\pi_j$.

For an ordered marker set $B$, set
\[
 P_j(f\epsilon_B)=
 \begin{cases}(\pi_jf)\epsilon_B,&j\in B,\\0,&j\notin B,
 \end{cases}
\]
\[
 h_j(f\epsilon_B)=
 \begin{cases}
 -(-1)^{\#\{i\in B:i<j\}}(S_jf)\epsilon_{B\cup\{j\}},&j\notin B,\\
 0,&j\in B.
 \end{cases}
\]
These maps are continuous coefficientwise.

Put
\[
 P=P_1\cdots P_r,
 \qquad h=\sum_{j=1}^rP_1\cdots P_{j-1}h_j.
\]
The image $\mathcal E_{r,s}$ of $P$ contains all normal
markers and only holomorphic normal coefficients. If
$\nu=\epsilon_1\cdots\epsilon_r$, then
\[
 \mathcal E_{r,s}^{-r-k}
 =\nu\,\mathcal J_r^{\mathrm{hol}}(W)
       \otimes\Lambda^k(\epsilon_{w_1},\ldots,\epsilon_{w_s}),
 \qquad d_E=(-1)^r d_w.
\]
Let $i$ be inclusion and $p$ projection, so $ip=P$.

\begin{proposition}\label{mht:prop:normal-source-contraction}
The maps satisfy
\[
 pi=1,\qquad dh+hd=1-ip,\qquad h^2=hi=ph=0.
\]
They commute with holomorphic source differential operators
and extension by zero in the tangent variables.
\end{proposition}
\begin{proof}
In one normal direction, integration gives
$dh_j+h_jd=1-P_j$ and $h_j^2=h_jP_j=P_jh_j=0$.
Its exterior sign makes $h_j$ anticommute with the
differential in every other direction. Each $P_j$ commutes
with all these operators. The telescoping sum gives
\[
 dh+hd=\sum_jP_1\cdots P_{j-1}(1-P_j)=1-P.
\]

In a product of two distinct homotopy summands, the later
one has a projection in the earlier normal direction.
This kills the earlier homotopy, in either order.
The side conditions follow.

Formal antiholomorphic integration and evaluation commute
with holomorphic multiplication and differentiation.
All maps preserve common tangent support and act on one
coefficient at a time. These facts prove the remaining
continuity and naturality statements.
\end{proof}

Write $A=\bar A j$ for the jet factor of an operation.
Give $\operatorname{Hom}_{\mathrm{cont}}(\mathcal E_{r,s},F)$
the Hom differential induced by $Q,d_E$.

\begin{theorem}\label{mht:thm:normal-operation-contraction}
There is a continuous cochain deformation retraction
\[
 \mathcal O_Z(F)\ \rightleftarrows\
                   \operatorname{Hom}_{\mathrm{cont}}(\mathcal E_{r,s},F)
\]
with maps
\[
 \mathbf P(A)=\bar A i,\qquad
 \mathbf I(B)=Bpj,\qquad
 \mathbf H(A)=(-1)^q\bar A h j\quad(|A|=q).
\]
It is natural in continuous cochain maps on $F$ and in
the individual holomorphic right differential actions.
If an operation has normal order at most $N$, projection
and inclusion preserve that bound and the homotopy lowers
it to $N-1$. The homotopy vanishes at order zero.
\end{theorem}
\begin{proof}
Taylor factorization identifies the supported space with
continuous Hom from the jet complex. Precomposition is
continuous on its bounded-set topology. For degree $q$,
\[
 \delta((-1)^qAh)=(-1)^qQAh+Ah d,
 \qquad
 (-1)^{q+1}(\delta A)h=(-1)^{q+1}QAh+Adh.
\]
Their sum is $A(1-ip)$. The contraction side conditions
follow by precomposition. The formulas commute with target
cochain maps and with holomorphic differentiation of tests.

Every term of $h$ raises normal monomial degree by one,
whereas $i,p$ never lower it. This proves the order bounds,
also when a bound is fixed only for one target seminorm.
\end{proof}

For one normal coordinate and scalar target, let $c_{ab}^0$
and $c_{ab}^1$ extract $[u^a\bar u^b]$ on the unmarked
and marked component. Then
\[
 \delta c_{ab}^0=(b+1)c_{a,b+1}^1,
 \qquad
 \mathbf Hc_{ab}^1=b^{-1}c_{a,b-1}^0\quad(b>0).
\]
The surviving classes $c_{a0}^1$ have degree one.
Their densities $\eta_a=(-1)^a\partial_u^a\delta_0/a!$
obey the right-module rules
\[
 \eta_a\partial_u=(a+1)\eta_{a+1},\qquad
 \eta_a u=\eta_{a-1}\ (a>0),\qquad \eta_0u=0.
\]
The positive sign in the right derivative comes from
precomposition on tests. It is distinct from the left
distributional derivative in \eqref{mht:eq:holomorphic-normal-actions}.

Formal completion permits increasing normal order.
For example, $\sum_{j\geq0}\chi^j c_{j0}^1$ is continuous
coefficientwise in $\chi$. No fixed finite sum of principal
parts represents it. Such a sum would annihilate every
sole coefficient $u^{N+1}$ beyond its order $N$, whereas
this operation returns $\chi^{N+1}$.
Section~\ref{mht:sec:adic-normal} distinguishes this
coefficientwise support from ordinary algebraic support.

Let $C_{\mathbf a}:\mathcal E_{r,s}\to\mathcal K(W)$ remove
$\nu$ and extract its $u^{\mathbf a}$ coefficient.
It has degree $r$ and satisfies
$C_{\mathbf a}d_E=(-1)^r d_w C_{\mathbf a}$.
For any continuous degree-zero cochain map
$M:\mathcal K(W)\to F$, the operation
\begin{equation}\label{mht:eq:supported-current-lift}
 A_{\mathbf a}(M)=MC_{\mathbf a}pj
\end{equation}
is a closed element of $\mathcal O_Z^r(F)$.
Indeed $QM=Md_w$ gives
$QA_{\mathbf a}=(-1)^rA_{\mathbf a}d_\partial$.
Coefficient differentiation and multiplication give its
normal right actions exactly as above.

An augmented target gives a simple nonzero class. Suppose
$1\in Z^0F$ and a continuous cochain map
$\varepsilon:F\to\mathbb C$ satisfies $\varepsilon(1)=1$.
Then
\[
 V_Z(\kappa\nu)=\pi^{-s}\int_W
                    \kappa_{\mathbf0,\mathbf0}(w)\,d^{2s}w\;1
\]
and zero on tangent-marker components defines a nonzero
class of degree $r$. Integration kills compact tangent
derivatives. After applying $\varepsilon$, the reduced
Hom complex has no degree below $r$, and $V_Z$ evaluates
nontrivially on a test of nonzero integral. It cannot be
a boundary.

When $Z$ has positive codimension, every separated annular
test vanishes near $Z$, so every supported operation kills
it. Such tests do not detect this class. The class belongs
to the occupied source set. The scalar of an empty
configuration is separate data.

Two adapted charts give two retractions of the same
supported Hom complex. Their transition is
$T_{ba}=\mathbf P_b\mathbf I_a$.
For three charts,
\[
 K_{cba}=-\mathbf P_c\mathbf H_b\mathbf I_a,
 \qquad
 d_cK_{cba}+K_{cba}d_a=T_{cb}T_{ba}-T_{ca}.
\]
This follows by inserting
$\delta\mathbf H_b+\mathbf H_b\delta=1-\mathbf I_b\mathbf P_b$.
Taking $c=a$ gives a homotopy inverse. These transitions
compare normal presentations on the same support, and
preserve the intrinsic holomorphic right actions.

\section{The tangent complex of the closed field}

The preceding operation model accepts any specified target
complex. It does not identify different closed theories.
This distinction already appears in tangent cohomology.

Put $\mathcal O=\mathbb C[[q,p]]$ and
$V=\mathcal O\otimes\Lambda(\theta_q,\theta_p)$, with
$|\theta_q|=|\theta_p|=1$ and $\pi=\theta_q\theta_p$.
Fix the signs
\begin{align*}
 \Delta(a\theta_q+b\theta_p)&=a_q+b_p,&
 \Delta(g\pi)&=g_q\theta_p-g_p\theta_q,\\
 d_\pi f&=f_p\theta_q-f_q\theta_p,&
 d_\pi(a\theta_q+b\theta_p)&=(a_q+b_p)\pi.
\end{align*}

The unspecified extreme components are zero. Direct
differentiation gives $d_\pi^2=\Delta^2=d_\pi\Delta+\Delta d_\pi=0$.
Define the fixed symplectic-background tangent complex
\[
 C_{\mathrm{symp}}=(V[[u]][2],d_\pi+u\Delta),
 \qquad |u|=2.
\]
The background is fixed before completion. This is not
evaluation at one of a background parameter completed
at zero.

For homogeneous coefficient degree $m$ define
\[
 Hf=0,\qquad
 H(a\theta_q+b\theta_p)=\frac{pa-qb}{m+1},\qquad
 H(g\pi)=\frac{qg\theta_q+pg\theta_p}{m+2}.
\]
Apply the formulas termwise to formal series and put
$p_u(f+a\theta_q+b\theta_p+g\pi)=f(0)-ug(0)$.
If $i$ includes constants, then
\begin{equation}\label{mht:eq:fixed-symplectic-contraction}
 (d_\pi+u\Delta)H+H(d_\pi+u\Delta)=1-ip_u.
\end{equation}

Indeed $\theta_q\mapsto dp$, $\theta_p\mapsto-dq$ identifies
the $d_\pi$ part with radial de Rham contraction.
On functions and vectors, $\Delta H+H\Delta=0$.
On a positive-degree bivector its two terms cancel, while
on a constant bivector it is evaluation at zero.
This gives exactly the term $-ug(0)$ in the projection.

For comparison, use formal mixed jets of the unmodified
geometric BF fields:
\[
 \mathcal A=\mathcal O[[r,s,\bar q,\bar p]]
       \otimes\Lambda(dr,ds,d\bar q,d\bar p),
\]
with the de Rham--Dolbeault differential in the four
displayed jet variables. Before a normal-support or
relative-source condition is imposed, its zero-field
tangent is
\[
 C_{\mathrm{BF}}=\mathcal A[[u]][1]\oplus\mathcal A[[u]][2],
\]
where $u$ is a spectator. Formal radial integration gives
a continuous contraction onto
$\mathcal O[[u]][1]\oplus\mathcal O[[u]][2]$ with zero differential.

\begin{proposition}\label{mht:prop:unchanged-tangent-obstruction}
The complexes $C_{\mathrm{BF}}$ and $C_{\mathrm{symp}}$
admit no quasi-isomorphism over $\mathbb C[[u]]$.
\end{proposition}
\begin{proof}
Equation~\eqref{mht:eq:fixed-symplectic-contraction} gives
$H^\bullet(C_{\mathrm{symp}})=\mathbb C[[u]][2]$.
Thus its degree $-1$ cohomology is zero. The BF contraction
gives $H^{-1}(C_{\mathrm{BF}})=\mathcal O$.
A quasi-isomorphism would identify these groups.
\end{proof}

A normally supported cone, another background, a quotient
or extra fields changes this comparison. Such a change
must be specified with its map. The proposition concerns
the two unchanged tangent complexes, and gives no
obstruction to every possible relative closed-field model.

\section{Closed operations with one normal current}

A brane source may fail the closed Maurer--Cartan equation.
Retaining that failure as curvature gives a precise local
model.

Work on a complex polydisk with coordinates $(p,q,z)$,
volume $dp\wedge dq\wedge dz$ and source divisor $D=\{p=0\}$.
Let $\xi_p,\xi_q,\xi_z$ be odd polyvector generators.
Antiholomorphic differentials precede them, and every
exterior generator has degree one. On smooth Dolbeault
polyvectors define
\[
 \bar\partial=\sum_xd\bar x\,\partial_{\bar x},\qquad
 \Delta=\sum_x\partial_x\partial^L_{\xi_x}.
\]
The left odd derivatives include signs from all preceding
exterior generators. The operators anticommute and square
to zero.

Let $\mathscr A$ be this smooth graded algebra and
$\mathscr M$ its module of normal distributions in $p$
with smooth distribution-valued dependence on $(q,z)$.
Smooth multiplication and differentiation act on $\mathscr M$.
Use the square-zero extension
\[
 \mathscr B=\mathscr A\ltimes\nu(\mathscr M),\qquad
 (a+\nu T)(b+\nu S)=ab+\nu(aS+Tb),
\]
with graded module signs. The symbol $\nu$ names a second
summand of degree zero. Its products with itself vanish.

In particular $\nu(1)$ is different from the algebra unit.
The identity $p\nu(1/p)=\nu(1)$ does not invert $p$ in
the unital algebra.

Complete in a source parameter $\lambda$ of degree zero
and a descendant $u$ of degree two. Set
\[
 E=\xi_p\xi_q,\qquad
 C=\nu(p^{-1}\xi_q),\qquad
 G=\nu(p^{-1}\xi_q\xi_z),\qquad
 K=\nu(\bar\partial(p^{-1})\xi_q),
\]
\[
 P=E+aG,\qquad a=\sigma\lambda,\qquad \sigma\in\{1,-1\}.
\]
The locally integrable function $1/p$ fixes the current.
With the positive complex orientation,
\begin{equation}\label{mht:eq:closed-source-residue}
 \bar\partial(1/p)=\pi\delta_0\,d\bar p,
 \qquad d(dp/p)=2\pi i\,\delta_0\,dx\wedge dy.
\end{equation}
The common $\xi_q$ gives $P^2=0$, and $\Delta P=0$.

On $\mathscr B[[\lambda,u]][2]$ use symmetric Taylor
operations. Before inserting $P$, let $b_0=0$,
$b_1=\bar\partial+u\Delta$ and, for $n\geq2$, set
\begin{equation}\label{mht:eq:closed-descendant-operations}
 b_n(u^{k_1}x_1,\ldots,u^{k_n}x_n)=
 \begin{cases}
 -\dfrac{(n-2)!}{k_1!\cdots k_n!}\Delta(x_1\cdots x_n),
                  &\sum_i k_i=n-2,\\[3pt]
 0,&\text{otherwise}.
 \end{cases}
\end{equation}
Each operation has degree one after the shift.
They are $\mathbb C[[\lambda]]$-multilinear, but are not
multilinear over the descendant ring.

\begin{theorem}\label{mht:thm:one-current-curved-operations}
Translation by $P$ defines a curved $L_\infty$ algebra.
Its curvature and unary operation are
\[
 b_0^P=aK\xi_z,
 \qquad
 b_1^Px=\bar\partial x+u\Delta x-\Delta(Px_0),
 \qquad x=\sum_{k\geq0}u^kx_k.
\]
For $n\geq2$, the translated $n$th operation consists of
\eqref{mht:eq:closed-descendant-operations} and the term
\[
 -\frac{(n-1)!}{k_1!\cdots k_n!}\Delta(Px_1\cdots x_n)
       \quad\text{when }\sum_i k_i=n-1.
\]
Restriction of the chart commutes with every operation.
An operation with one normal-distribution input preserves
its support, and an operation with two such inputs is zero.
\end{theorem}
\begin{proof}
For smooth compactly supported polyvectors, the classical
BCOV interaction has Taylor coefficients
\[
 I(\mu)=\sum_{m\geq3}\frac1{m!}
 \sum_{k_1+\cdots+k_m=m-3}
 \frac{(m-3)!}{k_1!\cdots k_m!}
              \operatorname{Tr}_\Omega(\mu_{k_1}\cdots\mu_{k_m}).
\]
Here $\operatorname{Tr}_\Omega(\gamma)$ is the integral of
$(\gamma\mathbin{\lrcorner}\Omega)\wedge\Omega$ and vanishes
outside top total degree. The genus-zero descendant recursion,
the second-order identity for $\Delta$ and integration by
parts give its classical master equation.\footnote{K.~Costello
and S.~Li, \emph{Quantum BCOV theory on Calabi--Yau manifolds
and the higher genus B-model},
\href{https://arxiv.org/abs/1201.4501}{arXiv:1201.4501}.
See Section~2.10, especially Lemma~2.10.2 and its proof.
The displayed state operators use the contragredient sign
convention for the nonlinear differential.}

Differentiating once and applying the primary Poisson kernel
gives \eqref{mht:eq:closed-descendant-operations}.
The state identities are the signed unshuffle sums
\[
 \sum_{I\sqcup J=[n]}\chi(I,J)
 b_{|J|+1}(b_{|I|}(x_I),x_J)=0.
\]

For one second-summand input, regularize its coefficient
in the normal variable. The desired identity is the term
linear in an auxiliary square-zero parameter in the smooth
identity. Each term contains derivatives of only one current
multiplied by smooth functions, so it passes to the
distributional limit. With two second-summand inputs every
nested term vanishes, either in the inner operation or
in the outer one.

Translation of a square-zero formal vector field preserves
its square. Here the Taylor translation is finite: all
products with two copies of $P$ vanish. Its constant term
is $\bar\partial P=aK\xi_z$, and its other terms give
the stated formulas. Smooth multiplication and differentiation
preserve distributional support and commute with restriction.
\end{proof}

The first curved identity is already nontrivial on a
holomorphic scalar $h(z)$:
\[
 b_1^Ph=ah'C,\qquad
 (b_1^P)^2h=ah'K,\qquad
 b_2^P(b_0^P,h)=-ah'K.
\]
The last equality uses the left odd derivative through
$d\bar p$. Thus the unary operator on the whole carrier
is not a differential for $a\ne0$.
The interaction itself extends to monomials with one
current by pairing it against compact smooth factors.
A nondegenerate BV pairing on the entire square-zero
extension, or a coupled disk action, is additional data.

There is nevertheless a supported subcomplex. Let
$\mathscr N\subset\mathscr M$ consist of currents supported
on $D$, locally of bounded finite normal order.
On $\nu(\mathscr N)[[\lambda,u]]$, the singular part of $P$
has zero product with every input. The unary operation is
\[
 d_E=\bar\partial+u\Delta-\Delta M_E\Pi_0,
 \qquad \Pi_0x=x_0,
\]
where $M_E$ is multiplication by $E$.
Define $S_ux=(x-x_0)/u$, $A_E=1-M_ES_u$ and
$Q_E=\bar\partial+[E,-]+u\Delta$, with
$[E,-]=M_E\Delta-\Delta M_E$.

\begin{proposition}\label{mht:prop:supported-closed-residue}
The continuous support-preserving conjugations are
\[
 d_E=A_EQ_EA_E^{-1},\qquad A_E^{-1}=1+M_ES_u,
\]
\[
 U_E=A_EuA_E^{-1}=u-M_E\Pi_0.
\]
The residue $K$ defines a nonzero class for $d_E$.
Its first transported descendant is exact:
\begin{equation}\label{mht:eq:closed-residue-descendant}
 U_EK=uK=d_E(zK\xi_z).
\end{equation}
\end{proposition}
\begin{proof}
The exterior factors give $M_E^2=M_E\Delta M_E=0$.
Also $\bar\partial E=0$, $S_uu=1$ and
$uS_u=1-\Pi_0$. Expanding the two conjugations proves
the formulas. Hence $U_E$ is a chain map. Ordinary
multiplication by $u$ need not be one, since
\[
 d_E(ux)-u d_Ex=u\Delta(Ex_0).
\]

To detect $K$, first quotient by positive descendant degree.
The differential becomes $\bar\partial-\Delta M_E$.
Next quotient by $\xi_p,\xi_z,d\bar q,d\bar z$.
This is a quotient complex: a term removing $\xi_z$
from $Ex$ still has the factor $E$, and $Ex=0$ if
$x$ already contains $\xi_p$.

Write $e_0=d\bar p$, $e_1=\xi_q$ and represent the
distribution $\partial_p^i\partial_{\bar p}^j\delta_0$
by $R^iT^j$. The quotient is
\[
 \bigl(C^\infty(D)[R,T]\otimes\Lambda(e_0,e_1)\bigr)[[\lambda]],
 \qquad d=T e_0-R e_1.
\]
The sequence $R,T$ is regular, so its degree-two cohomology
is $C^\infty(D)[[\lambda]]e_0e_1$. The image of $K$ is
$\pi e_0e_1$, detected by setting $R=T=0$.
Thus $K$ was not exact in the original complex.

Finally $\bar\partial K=\Delta K=EK=0$ and
$\Delta(zK\xi_z)=K$. The other terms of $d_E$ on
$zK\xi_z$ vanish, proving \eqref{mht:eq:closed-residue-descendant}.
\end{proof}

The normal orientation retains a boundary term. Let $\chi$
be a radial cutoff equal to one near $p=0$ and zero near
the outer circle. For a compact tangential top form $\psi_D$
of integral one,
\[
 \int d\chi\wedge\frac{dp}{p}\wedge\psi_D=-2\pi i.
\]
The inner boundary has the opposite of the positive-circle
orientation. Equation~\eqref{mht:eq:closed-source-residue}
supplies the corresponding positive supported term.
Removing the normal divisor would discard this information.

\section{Two sources and collision extensions}

A unit one-source residue does not determine a two-source
contact. Write $\delta_a(z)=e^{-|z|^2/a}/(\pi a)$ and let
$\Delta_{\mathbb R^2}=\partial_x^2+\partial_y^2$.
For $\alpha,\beta>0$, multiplication of the smooth densities gives
\[
 \delta_{\alpha s}\delta_{\beta s}
 =\frac1{\pi(\alpha+\beta)s}
             \delta_{\alpha\beta s/(\alpha+\beta)}.
\]

The centered Gaussian has zero odd moments and second
moments $a/2$ in each real coordinate. Taylor expansion
through degree three therefore yields
\begin{equation}\label{mht:eq:two-source-width-contact}
 \delta_{\alpha s}\delta_{\beta s}
 =\frac{\delta_0}{\pi(\alpha+\beta)s}
 +\frac{\alpha\beta}{4\pi(\alpha+\beta)^2}
                         \Delta_{\mathbb R^2}\delta_0+O(s).
\end{equation}
The remainder on a compact smooth test is bounded by a
constant times $s$ and its fourth-derivative seminorm.

The finite coefficient is $1/(16\pi)$ for equal widths
and $1/(18\pi)$ for width ratio two. These are scalar
coefficients after contraction. The calculation does not
multiply two identical Dolbeault one-forms.

More generally, let $t$ be a distribution on
$\mathbb R^d\setminus\{0\}$ of finite scaling degree $s$.
This means that $\lambda^{s'}t(\lambda\,\cdot)$ tends to
zero distributionally for every $s'>s$.
For $s\geq d$ put $m=\lfloor s-d\rfloor$ and choose
$\chi=1$ near zero with compact support. Taylor subtraction is
\begin{equation}\label{mht:eq:collision-taylor-subtraction}
 R_\chi t(\varphi)=t\left(\varphi-\chi
       \sum_{|\nu|\leq m}\frac{x^\nu}{\nu!}
                                  \partial^\nu\varphi(0)\right).
\end{equation}

The pairing is the limit after removing a shrinking
neighborhood of zero. On dyadic annuli, the Taylor remainder
gives the convergent bound $2^{-j(d+m+1-s')}$ for
$s<s'<d+m+1$. This proves existence and continuity.
The rescaled estimate gives scaling degree at most $s$.
For $s<d$ the same annular estimate requires no subtraction.

Two extensions with this scaling bound differ by
derivatives of delta of order at most $m$.
Indeed their difference is supported at zero, and a
distribution supported at one point has finite derivative
order. Its highest derivative has scaling degree
$d+|\nu|$, which gives the bound. In particular
\[
 (R_{\chi'}-R_\chi)t(\varphi)
 =-\sum_{|\nu|\leq m}\frac{\partial^\nu\varphi(0)}{\nu!}
                          t((\chi'-\chi)x^\nu).
\]
For $s<d$ the extension with the stated bound is unique.
Nested and overlapping collisions still require compatible
choices on the full collision diagram.

Closedness adds an independent condition. Let
$Z_0\subset Z_1$ be closed source sets and put
$V^\circ=V\setminus Z_0$. Suppose $A^\circ$ is a closed
degree-$q$ operation supported on $Z_1\cap V^\circ$.

\begin{proposition}\label{mht:prop:supported-extension-obstruction}
If $A^\circ$ has a continuous degree-$q$ extension
$E\in\mathcal O_{Z_1}(F)$, then
\[
 [\delta E]\in H^{q+1}\mathcal O_{Z_0}(F)
\]
is independent of the extension. It vanishes exactly
when there is a closed extension supported on $Z_1$.
For a linear partial diagonal, its reduced representative
is $\mathbf P(\delta E)$.
\end{proposition}
\begin{proof}
Restriction is precomposition by extension of compact
tests by zero and commutes with $\delta$.
Thus $\delta E$ is supported on $Z_0$ and is closed.
Two extensions differ by an operation supported on $Z_0$,
so their differentials differ by a boundary there.

If $\delta E=\delta C$ for a supported degree-$q$
operation $C$, then $E-C$ is a closed extension.
Conversely, a closed extension subtracted from $E$ gives
such a $C$. The normal-operation contraction preserves
the class and supplies the reduced representative.
\end{proof}

This cohomology uses the actual image of $\delta$, not
its topological closure. If the permitted local corrections
form a subcomplex $\mathcal L\subset\mathcal O_{Z_0}(F)$,
the chosen $E$ can be corrected only when
$\delta E\in\delta(\mathcal L^q)$.
Vanishing in the larger complex need not imply this.
Existence of a continuous extension is also a hypothesis.

For $n$ occupied positions, shifting the operation complex
by $[n-1]$ puts the current of an $s$-block partition in
degree $r-(n-1)=1-s$. A three-position pair-diagonal
current thus has degree $-1$, while a small-diagonal
current has degree zero. This permits a homotopy equation
between them. The degree calculation does not construct
the homotopy or identify its boundary with a collision defect.

\section{A four-position attachment}

Even a scalar collision coefficient can have several
cohomology modules joined by a nonzero extension.
Take $X=\mathbb A^1_{\mathbb C}$ and four labeled positions.
Let
\[
 D=\Delta_{14},\qquad P=\Delta_{13}\cap\Delta_{24},
 \qquad \Delta=D\cap P,
 \qquad Z=D\cup P.
\]
Here $D$ is a divisor, $P$ a codimension-two surface,
and $\Delta$ the small diagonal. Work with right algebraic
differential modules and define the specific coefficient
\[
 \mathcal K_4=R\Gamma_Z(\omega_{X^4})[3].
\]

This is a local-cohomology construction. Identifying it
with another coefficient system requires a comparison map.

\begin{proposition}\label{mht:prop:four-position-attachment}
The only cohomology modules of $\mathcal K_4$ are
\[
 H^{-2}(\mathcal K_4)=i_{D,*}\omega_D,
 \qquad
 H^{-1}(\mathcal K_4)=i_{P,*}\omega_P(*\Delta).
\]
The resulting extension is nonzero. Its class spans
\[
 \operatorname{Ext}^2_{\mathcal D_{X^4}}
       (i_{P,*}\omega_P(*\Delta),i_{D,*}\omega_D)
       \cong\mathbb C.
\]
It is the principal part along $\Delta$ in $P$, followed
by the codimension-two local-cohomology boundary in $D$.
\end{proposition}
\begin{proof}
Use coordinates $a=x_1-x_4$, $b=x_1-x_3$,
$c=x_2-x_4$, $w=x_4$. Then $D=(a)$, $P=(b,c)$
and $\Delta=(a,b,c)$. The localization triangle for the
pair $D\subset Z$ is
\[
 R\Gamma_D\omega\longrightarrow R\Gamma_Z\omega
 \longrightarrow R\Gamma_P(\omega[a^{-1}])
 \longrightarrow R\Gamma_D\omega[1].
\]
The first term has cohomology in degree one and the third
in degree two. Their regular normal sequences give the
two displayed modules after the shift by three.

To compute the connecting class, tensor with the inverse
density and use left differential modules in these coordinates.
In one variable, write $O=\mathbb C[x]$ and
$M=\mathbb C[x,x^{-1}]/\mathbb C[x]$.
The two-term differential-module resolution of $M$ by
right multiplication by $x$ gives
$\operatorname{Hom}_{\mathcal D}(M,O)=0$ and
$\operatorname{Ext}^1_{\mathcal D}(M,O)=\mathbb C$.
The latter is detected by evaluation modulo $x$.

Also $\operatorname{Hom}_{\mathcal D}(O[x^{-1}],M)$
is a line generated by the principal part. A map from
$O$ to $M$ must send $1$ to a vector killed by
$\partial_x$. Differentiation is injective on $M$, so
that map is zero. Every map therefore factors through
the principal-part quotient, whose differential-module
endomorphisms are scalars.

In the $a$ direction the connecting map is this principal
part. In the $b,c$ directions it is the tensor product
of the two nonzero degree-one extension classes.
The $w$ direction has the identity of $O_w$.
The tensor-product resolutions therefore give the stated
one-dimensional Ext group. The localization boundary
has coefficient one in each factor, so its class is nonzero.
A direct sum of the shifted cohomology modules would
have zero connecting class.
\end{proof}

The crossing module has poles along the small diagonal.
Keeping only regular coefficients on $P$ loses its
principal-part boundary. Keeping the two cohomology modules
but discarding their attachment also changes $\mathcal K_4$.
A supported operation comparison must preserve both pieces
and the connecting map.

\section{A relative polar source algebra}

The square-zero closed carrier and principal parts have
different source products. There is a simple reason not
to multiply all normal currents in an ordinary unital algebra.
If a central element $z$ has an inverse $c$ with $zc=1$
and annihilates $\delta$, then
\[
 \delta=(cz)\delta=c(z\delta)=0.
\]
Thus a nonzero source satisfying $z\delta=0$ cannot
coexist with that central inverse. A relative operator
algebra retains the source by making the inverse one-sided.

On $\mathcal P_z$ write $e_n=[z^{-n-1}]$ and define
\[
 Le_0=0,\qquad Le_n=e_{n-1}\ (n>0),\qquad Se_n=e_{n+1}.
\]
Let $P_0$ project onto $e_0$. Then $LS=1$ and
$SL=1-P_0$. The algebra $\mathcal T$ generated by $L,S$
has finite-matrix ideal
\[
 \mathcal F=\bigoplus_{i,j\geq0}\mathbb C E_{ij},
 \qquad E_{ij}=S^iP_0L^j,
 \qquad E_{ij}E_{kl}=\delta_{jk}E_{il}.
\]
Here finite matrices are matrices with finitely many
nonzero entries in this basis. They are not all
finite-rank endomorphisms of an algebraic direct sum.

The symbol sequence is exact:
\begin{equation}\label{mht:eq:polar-symbol-extension}
 0\longrightarrow\mathcal F\longrightarrow\mathcal T
 \xrightarrow{\sigma}\mathbb C[\zeta,\zeta^{-1}]
 \longrightarrow0,
 \qquad \sigma(S)=\zeta,\quad\sigma(L)=\zeta^{-1}.
\end{equation}
Indeed $LS=1$ reduces every word to $S^iL^j$.
Repeated use of $SL=1-P_0$ expresses it as a Laurent
lift plus a finite matrix. A nonzero Laurent symbol
acts on all sufficiently high basis vectors and cannot
be a finite matrix. This proves exactness and faithfulness
of the shift presentation.

For $N\geq1$, the projector
\[
 \Pi_N=1-S^NL^N=[L^N,S^N]=\sum_{i=0}^{N-1}E_{ii}
\]
has corner $\Pi_N\mathcal T\Pi_N\cong\operatorname{Mat}_N$,
with corner unit $\Pi_N$ and trace $N$.
The trace pairing is perfect since $E_{ij}$ pairs with
$E_{ji}$. At rank two,
$\operatorname{Tr}([E_{01},E_{10}]^2)=2$.
Diagonal projection kills both inserted matrix units and
gives zero, despite preserving the trace of the unit.

Choose $T_{\zeta^a}=S^a$ for $a\geq0$ and
$T_{\zeta^{-b}}=L^b$ for $b>0$. The contact is
$c(f,g)=T_fT_g-T_{fg}\in\mathcal F$.
For $a,b>0$,
\begin{equation}\label{mht:eq:polar-section-contact}
 c(\zeta^a,\zeta^{-b})
   =-\sum_{j=\max(0,b-a)}^{b-1}E_{j+a-b,j}.
\end{equation}
The product $S^aL^b$ kills $e_j$ for $j<b$.
Its Laurent lift differs exactly at the indices in the
sum, proving the formula. Reverse-order and same-sign
monomial contacts vanish.

Associativity gives the full
factor identity
\[
 T_fc(g,h)-c(fg,h)+c(f,gh)-c(f,g)T_h=0.
\]
For example, $c(\zeta^3,\zeta^{-2})=-E_{10}-E_{21}$.

The finite-matrix trace gives
\begin{equation}\label{mht:eq:polar-residue-transgression}
 \operatorname{Tr}_{\mathcal F}[T_f,T_g]
       =\operatorname{Res}_{\zeta=0}f\,dg=:\omega(f,g).
\end{equation}
For monomials $f=\zeta^a$, $g=\zeta^b$, a nonzero
trace requires $a+b=0$ and has value $b$.
These are exactly the residue coefficients.
Changing a lift by a finite matrix changes its commutator
by terms of trace zero. Antisymmetry and the product
rule show that $\omega$ is a cyclic one-cocycle.
No cyclic trace on all of $\mathcal T$ can extend this
trace, since $[L,S]=P_0$ has trace one.

Every element is uniquely $T_f+F$. Define
$\widetilde\tau(T_f+F)=\operatorname{Tr}_{\mathcal F}F$.
With the bare scalar Hochschild differential,
\[
 b\widetilde\tau=\sigma^*\omega.
\]
The finite-matrix cross commutators have zero trace,
and \eqref{mht:eq:polar-residue-transgression} gives the
remaining term. Consequently $(\widetilde\tau,\omega)$
is closed in the relative cyclic complex with differential
\[
 d(\tau,\gamma)=(b\tau-\sigma^*\gamma,-b\gamma).
\]
Its degree-$n$ term is
$C^n_\lambda(\mathcal T)\oplus C^{n+1}_\lambda
(\mathbb C[\zeta,\zeta^{-1}])$, where cyclic cochains
obey $\phi(a_n,a_0,\ldots,a_{n-1})=(-1)^n\phi(a_0,\ldots,a_n)$.

Normal derivatives also act on $\mathcal P_z$:
$\partial_z e_n=-(n+1)e_{n+1}$ and $[\partial_z,L]=1$.
Adjoining this operator enlarges $\mathcal T$.
Independent normal directions use tensor products, and
their contacts compose as operators. Since $P_0^2=P_0$,
no product-preserving map can send a square-zero source
to $P_0$. The native action must determine which source
product its collisions realize.

The one-sided inverse also gives a homotopy obstruction.
Let $D$ be a unital dg algebra with closed $S,L$ in
degree zero, $LS=1$, $SL=1-e$. Suppose $eDe$ has a
closed degree-zero functional $\tau$ with $\tau(e)=1$.
Require the cohomology images $\zeta\mapsto S$, $\zeta^{-1}\mapsto L$.
There is no such cochain map from the Laurent algebra to $D$
that is unital and multiplicative up to homotopy.

To prove this, write
$F_1(1)=1+dh_0$, $F_1(\zeta)=S+dx$ and
$F_1(\zeta^{-1})=L+dy$.
If $dH(a,b)=F_1(a)F_1(b)-F_1(ab)$, then
\[
 h=-H(\zeta,\zeta^{-1})-h_0+xL+Sy+x\,dy
 \quad\text{satisfies}\quad dh=e.
\]
The closed idempotent gives $d(ehe)=e$ inside the corner.
Applying $\tau$ would give $1=0$. This obstruction retains
the prescribed generators and corner detector. It does
not exclude a relative comparison such as the cyclic
pair above.

\section{The joint source and quantum comparison}

The coupled classical theory admits more loops than its
pure cotangent sector. In a connected graph, let $v_B$
be the number of bulk vertices, $v_C$ coupling vertices
and $v_O$ open cubic vertices. If $b$ cotangent and $a$
open legs are external, then
\[
 e_B=v_B-b,\qquad e_A=\frac{2v_C+3v_O-a}{2},\qquad
 b_1=1-b+\frac{v_O-a}{2}.
\]
Two open cubic vertices joined by three propagators
already have two loops. Any vanishing of their weight
requires a geometric calculation.

In a finite BV algebra, or on a specified renormalized
domain where every term is defined, write
$I=I_0+\hbar I_1+\hbar^2I_2+\cdots$.
Suppose $I_0$ solves the classical master equation and
put $Q_{I_0}=Q+\{I_0,-\}$. The quantum equation gives
\begin{align}
 Q_{I_0}I_1&=-\Delta I_0,\notag\\
 Q_{I_0}I_n&=-\Delta I_{n-1}
       -\tfrac12\sum_{a=1}^{n-1}\{I_a,I_{n-a}\},\qquad n\geq2.
       \label{mht:eq:coupled-supported-QME}
\end{align}
The BV identities make the next right side closed once
the preceding equations hold. A solution requires its
primitive in the permitted local cyclic complex, including
its normal support and boundary conditions.

If a finite regularization realizes the change of variables
in \eqref{mht:eq:surviving-cyclic-cubic} canonically, its
triangular matrix has identity diagonal and Berezinian one.
It preserves the finite BV measure and Laplacian.
Quantum master solutions for the separated open sector
and the bulk sector with its residual cubic therefore
pull back to a coupled solution. This implication requires
both separated solutions. It does not construct a
continuum regulator or its allowed counterterms.

A source comparison must retain its homotopy as well.
For cochain maps $f:L_{\mathrm{nat}}\to L$ and
$g:L_{\mathrm{cyc}}\to L$, the joint tangent complex is
\[
 \mathcal P^n=L_{\mathrm{nat}}^n\oplus L_{\mathrm{cyc}}^n
                      \oplus L^{n-1},
 \qquad
 D(a,b,h)=(da,db,f(a)-g(b)-dh).
\]
The two cochain-map identities give $D^2=0$.
The component $h$ is the comparison homotopy.

For nonlinear deformation theory, use the path-object
homotopy pullback of the dg Lie algebras. The linear
cone alone supplies no Lie bracket. In the deformation
complex, the next Maurer--Cartan obstruction has degree two.

\begin{problem}\label{mht:qn:native-supported-source}
Fix the polar action, its Neumann boundary domain and
its normal lattice. Construct a map from independently
defined sourced closed fields to continuous cyclic
operations of the open fields. Determine the
permitted normal markings and structured homotopies from
these physical data. Match the quadratic pair, first
ghost descendants, unit, support and cyclic pairing.

Extend the map to simultaneous source insertions and
all supported collision and sewing operations. Preserve
the four-position attachment and identify the image of
the allowed local counterterms in the supported operation
complex. For a quantum comparison, construct a common
regulator and solve \eqref{mht:eq:coupled-supported-QME}.
\end{problem}

The real internal defect, the holomorphic normal source
and a spectral brane have independently specified supports.
The cyclic cap and the scalar relative models give
constructions on their stated domains. Their identification
with the full sourced closed multiplet is the additional
map required by the problem.

\section{Completion in one normal direction}
\label{mht:sec:adic-normal}

Increasing normal order is compatible with formal coefficients,
but changes what it means for a coefficient to be supported.
We isolate this issue in one algebraic normal direction.
The resulting reconstruction and convolution have explicit
maps. Their comparison with the sourced fields of
Problem~\ref{mht:qn:native-supported-source} remains separate.

Put $R=\mathbb C[[\hbar]]$ and $R_r=R/(\hbar^r)$.
Here $\hbar$ is a coefficient parameter independent of
the normal coordinate $u$. Let
\[
 D_R=R\langle u,\partial\rangle/(\partial u-u\partial-1),
 \qquad P_R=R[u,u^{-1}]/R[u].
\]

We use right differential modules. In the density basis
$\beta_m=[u^{-m-1}]$, their normal actions are
\[
 \beta_0u=0,\qquad \beta_mu=\beta_{m-1}\quad(m>0),
 \qquad \beta_m\partial=(m+1)\beta_{m+1}.
\]
Thus $P_R=D_R/uD_R$, with a free right-module resolution
\begin{equation}\label{mht:eq:normal-free-resolution}
 [D_R\xrightarrow{L_u}D_R]_{[-1,0]}\longrightarrow P_R.
\end{equation}
The map $L_u$ is left multiplication by $u$.
The shift convention is $X[s]^i=X^{i+s}$.

Work in the derived category of right $D_R$-modules.
For a complex $X$, set
\[
 X_r=X\otimes_R^{\mathbb L}R_r,
 \qquad \Lambda X=\mathbf R\varprojlim_r X_r.
\]
Ordinary open restriction $X[u^{-1}]$ localizes its underlying
$R[u]$-complex and retains the induced differential action.
It is exact.

We call $X$ \emph{derived adically supported}
if $X_r[u^{-1}]\simeq0$ for every $r$.
Let $\mathcal C_\hbar$ be the full category of such complexes
for which $X\to\Lambda X$ is an equivalence, and let
$\mathcal G_\hbar$ be the category of derived-complete
$R$-complexes. Neither category has a boundedness condition.
There is no completion in $u$.

\begin{lemma}\label{mht:lem:normal-completion-precision}
The completion map induces
$(\Lambda X)_r\simeq X_r$ for every complex $X$ and every $r$.
Moreover,
\[
 X_r[u^{-1}]\simeq0\ \text{for all }r
 \quad\Longleftrightarrow\quad X_1[u^{-1}]\simeq0
 \quad\Longleftrightarrow\quad\Lambda(X[u^{-1}])\simeq0.
\]
Completion preserves derived adic support.
\end{lemma}
\begin{proof}
The perfect scalar complex $[R\xrightarrow{\hbar^r}R]$
computes $R_r$, so its tensor commutes with homotopy limits.
Put $Y=X_r$. For $n\geq r$, the resolution computing
$Y\otimes_R^{\mathbb L}R_n$ has zero differential in its
two-term resolution direction. Its transition is the identity
on the degree-zero term and multiplication by $\hbar$
on the degree-minus-one term.

The second tower has zero composite after $r$ transitions.
In the product model for its homotopy limit, $1-\mathrm{shift}$
has the finite geometric-series inverse. That limit vanishes,
while the first tower has limit $Y$. This proves the assertion
about the natural completion map, also for unbounded $X$.

Put $N=X[u^{-1}]$. Exact localization commutes with
tensoring by $R_r$. The condition $N_1\simeq0$ says that multiplication
by $\hbar$ is an equivalence on $N$. It is equivalent to
$N_r\simeq0$ for all $r$, and hence to $\Lambda N\simeq0$
by the first assertion. Applying that assertion again proves
preservation of support.
\end{proof}

For an $R$-flat module, the support condition is equivalent to
\[
 \text{for every }v\text{ and }r\text{ there is }N
             \text{ such that }vu^N\in\hbar^rX.
\]
For a general complex the condition concerns all cohomology
of the derived reductions, including their Tor terms.
Ordinary quotient modules alone do not supply that condition.

\begin{example}\label{mht:ex:completed-principal-parts}
The completion of $P_R$ is the flat module
\[
 \widehat P_R=
 \left\{\sum_{m\geq0}a_m\beta_m:
       a_m\in R,\ a_m\longrightarrow0\ \hbar\text{-adically}\right\}.
\]
Its reduction is $P_{R_r}$, so it lies in $\mathcal C_\hbar$.
Nevertheless it has nonzero ordinary open restriction.
Indeed
\[
 p=\sum_{m\geq0}\hbar^m\beta_m
 \qquad\text{satisfies}\qquad pu=\hbar p,
\]
and no power of $u$ kills $p$.

More precisely, its $u$-torsion is exactly $P_R$.
Lowering is surjective, so
$\widehat P_R[u^{-1}]=\widehat P_R/P_R$.
Multiplication by $\hbar$ is invertible on this quotient.
Subtract a finite representative modulo $\hbar$ and divide
to prove surjectivity. Finite support proves injectivity.
The open quotient is therefore nonzero but
has zero derived completion.

Any complete complex with the specified compatible reductions
$P_{R_r}$ is their derived inverse limit $\widehat P_R$.
Consequently this tower has no complete realization with
ordinary $u$-support. Its derived adic realization exists.
\end{example}

Completion is left adjoint to the inclusion of complete
complexes. Taking the inverse limit of the multiplication
triangles gives
\[
 T(X)\longrightarrow X\longrightarrow\Lambda X,
 \qquad T(X)=\mathbf R\varprojlim(\cdots\xrightarrow{\hbar}X
                                      \xrightarrow{\hbar}X).
\]
Multiplication by $\hbar$ is invertible on $T(X)$.
Maps from such a complex to a complete complex vanish:
after reduction the same scalar action is both invertible
and nilpotent on the mapping complex. This proves the
stated adjunction.

\subsection{Reconstruction from the normal ground}

Define induction and restriction by
\[
 I(E)=\Lambda(P_R\otimes_R^{\mathbb L}E),
 \qquad K(M)=\operatorname{RHom}_{D_R}(P_R,M).
\]
Resolution~\eqref{mht:eq:normal-free-resolution} computes
$K(M)$ as $[M\xrightarrow{\cdot u}M]$ in normal degrees
$0,1$. This arrow is $R$-linear, without right
$D_R$-linearity.

\begin{theorem}\label{mht:thm:adic-normal-reconstruction}
The adjunction maps give inverse equivalences
\[
 I:\mathcal G_\hbar\ \rightleftarrows\
                       \mathcal C_\hbar:K.
\]
No flatness or finite-rank hypothesis on the scalar complex
is required.
\end{theorem}
\begin{proof}
Let $B$ be a commutative complex algebra. Write $D_B,P_B$
for the normal algebra and module over $B$, and put
$K_B(V)=\operatorname{RHom}_{D_B}(P_B,V)$.
First suppose $V$ is a module with locally nilpotent $u$-action.
The finite Taylor projector is
\[
 \pi(v)=\sum_{j\geq0}\frac{(-1)^j}{j!}(vu^j)\partial^j.
\]
The identity $\partial^ju=u\partial^j+j\partial^{j-1}$
gives $\pi(v)u=0$. With $E=\ker u$, define
\[
 P_B\otimes_B E\longrightarrow V,
 \qquad \beta_m\otimes e\longmapsto e\partial^m/m!.
\]

It respects both right actions. Applying the highest power
of $u$ in a finite relation proves injectivity.
Surjectivity follows from
\[
 v=\sum_{m\geq0}\frac{\pi(vu^m)\partial^m}{m!}.
\]
On expansion, the coefficient in total degree $l$ is
$\sum_{j=0}^l(-1)^j/(j!(l-j)!)$, equal to zero for $l>0$
and one for $l=0$. In particular, $u:V\to V$ is surjective.

Now suppose the cohomology of a complex $V$ is locally
$u$-nilpotent. The two-term restriction complex gives
$H^qK_B(V)=\ker(u:H^qV\to H^qV)$, because the preceding
cokernel vanishes. Since $P_B$ is free over $B$, the
preceding module calculation proves that evaluation
$P_B\otimes_B^{\mathbb L}K_B(V)\to V$ is an equivalence.
The unit is an equivalence because $u:P_B\to P_B$ is a
split scalar surjection with kernel $B\beta_0$.

Apply this argument over $B=R_r$. The finite free resolution
gives $K(M)_r\simeq K_r(M_r)$ and
$K(\Lambda M)\simeq\Lambda K(M)$.
Thus $K$ preserves completeness, and the reductions of
$I(E)$ are $P_{R_r}\otimes_{R_r}^{\mathbb L}E_r$.
The unit and counit are equivalences at every precision.
Both sides are complete, so they are equivalences before
reduction as well.
\end{proof}

For a complete differential-module complex $X$, evaluation
$IK(X)\to X$ is consequently right adjoint to the inclusion
of $\mathcal C_\hbar$. Limits taken in this supported
category are
\[
 \varprojlim_{\mathcal C_\hbar}M_i
       \simeq I\bigl(\mathbf R\varprojlim_i K(M_i)\bigr).
\]
They need not agree with ambient differential-module limits.

\subsection{Additive convolution and relative tensors}

The normal carrier has a geometric operation defined before
using the equivalence. Let $a(x,y)=x+y=t$ and let $a_+$ be
right differential-module direct image, computed by the
relative de Rham complex. Set
\[
 M*_+N=\Lambda a_+(M\boxtimes_R^{\mathbb L}N),
 \qquad E\widehat\otimes_R F=\Lambda(E\otimes_R^{\mathbb L}F).
\]
The second operation makes $\mathcal G_\hbar$ symmetric
monoidal: the kernel of completion consists of the
$\hbar$-invertible complexes and is a tensor ideal.

\begin{proposition}\label{mht:prop:adic-normal-convolution}
There is a natural equivalence
\[
 I(E)*_+I(F)\simeq I(E\widehat\otimes_R F).
\]
It makes $I$ a symmetric monoidal equivalence, with
convolution unit $\widehat P_R$.
\end{proposition}
\begin{proof}
At precision $r$, change coordinates to $t=x+y$, $s=y$.
This identifies the normal distribution module
$P_{x,R_r}\boxtimes P_{y,R_r}$ with
$P_{t,R_r}\boxtimes P_{s,R_r}$.
The relative de Rham complex is
\[
 [P_{s,R_r}\xrightarrow{\partial_s}P_{s,R_r}]_{[-1,0]}.
\]
Its differential sends $\beta_m$ to $(m+1)\beta_{m+1}$.
It is injective, with free cokernel $R_r\beta_0$ and a
scalar splitting. Hence the direct image is $P_{t,R_r}$,
with no shift, also after tensoring with arbitrary scalar
complexes.

The coordinate change sends
\[
 \beta_i(x)\beta_j(y)\longmapsto
       \sum_{k=0}^j\binom{i+k}{k}
                       \beta_{i+k}(t)\beta_{j-k}(s).
\]
Relative de Rham projection therefore gives
$\binom{i+j}{i}\beta_{i+j}(t)$.
The multinomial identity gives the same result for three
inputs in either order. Addition and iterated relative
de Rham complexes supply the coherent associativity and
symmetry maps. The origin inclusion supplies the unit.

Scalar reduction commutes with these complexes.
Lemma~\ref{mht:lem:normal-completion-precision} shows that
completing the inputs changes none of their reductions.
The comparison is an equivalence at each precision.
Completion proves the assertion.
\end{proof}

Let $B$ be an associative algebra in $\mathcal G_\hbar$,
with a right module $E$ and a left module $F$.
Transporting their actual actions through the equivalence
and realizing the two-sided bar gives
\begin{equation}\label{mht:eq:normal-relative-convolution}
 I(E)*_{I(B)}^{\mathbb L}I(F)
       \simeq I\bigl(\Lambda(E\otimes_B^{\mathbb L}F)\bigr).
\end{equation}
Indeed the bar terms correspond under the monoidal
equivalence, including every face and degeneracy.
Colimits in the complete scalar category are completed
ambient colimits. This proves the formula and support
of the result without a flatness assumption.

Flatness itself is not preserved. Take
$B=\widehat{R[t]}^{\,(\hbar)}$, the series with coefficients
tending to zero, and the two complete flat character
modules $E_0=R$, $E_\hbar=R$ with $t=0,\hbar$.
The free resolution of $E_0$ has differential $t$, so
\[
 E_0\otimes_B^{\mathbb L}E_\hbar
       =[R\xrightarrow{\hbar}R]_{[-1,0]}
       \simeq R/(\hbar).
\]
At precision $r$, this complex has both
$H^{-1}=\mathbb C\hbar^{r-1}$ and $H^0=\mathbb C$.
Both remain supported after normal induction.
A flat-only coefficient class would lose this relative tensor.

Two other operations have different output types.
Write $(-)^t$ for transpose under
$u\mapsto u$, $\partial\mapsto-\partial$. Then
\[
 \Lambda\bigl(I(E)\otimes_{D_R}^{\mathbb L}I(F)^t\bigr)
       \simeq(E\widehat\otimes_R F)[1],
 \qquad
 \operatorname{RHom}_{D_R}(I(E),I(F))
       \simeq\operatorname{RHom}_R(E,F).
\]
For the first formula, the normal resolution gives the
surjective map $u:P_{R_r}^t\to P_{R_r}^t$ in degrees
$-1,0$, with kernel $R_r$. The second is full faithfulness.

Both outputs are scalar complexes. Neither retains a
normal coordinate on which to impose support.

\subsection{Products, duality and the class of tests}

\begin{proposition}\label{mht:prop:normal-product-failure}
The ambient product $W=\prod_{i\geq0}\widehat P_R$ is
flat and complete, but is not derived adically supported.
\end{proposition}
\begin{proof}
It is torsion-free over the discrete valuation ring $R$,
hence flat. Componentwise division and inverse limits give
$W/\hbar^rW=\prod_iP_{R_r}$ and completeness.
In its special fiber, let $e=(\beta_i)_{i\geq0}$.
For every $N$, coordinate $i=N$ of $eu^N$ is $\beta_0$.
Thus no power of $u$ kills $e$, and $W_1[u^{-1}]\ne0$.
\end{proof}

The product inside $\mathcal C_\hbar$ is instead
$U=I(R^{\mathbb N})$. In coefficients $a_{i,m}$ its
inclusion into $W$ imposes
\[
 \text{for every }r\text{ there is }M_r\text{ such that }
 a_{i,m}\in\hbar^rR\quad\text{for all }i\text{ and }m\geq M_r.
\]
The ambient product permits the mode bound to depend on $i$.
The diagonal family above detects the nonzero quotient $L=W/U$.

This inclusion is $\hbar$-saturated, so $L$ is flat. As a
cofiber of complete complexes it is derived complete.
Lowering is surjective on $U$ and $W$, with the same kernel
$R^{\mathbb N}\beta_0$. Hence $u$ is invertible on $L$ and
\[
 K(L)=0,\qquad\operatorname{RHom}_{D_R}(D_R,L)=L\ne0.
\]
The normal ground misses a quotient that the free
differential-module test detects.

Define the ordinary normalized coefficient dual with
complete target by
\[
 \mathbb D_\hbar(M)=
       \operatorname{RHom}_{D_R}(M,\widehat D_R)^t[1],
 \qquad \widehat D_R=\Lambda D_R.
\]
The shift is the dimension-one normalization.
This dual takes values in ordinary differential modules.

\begin{proposition}\label{mht:prop:normal-realized-dual}
The flat supported complete module
$M=\Lambda(\bigoplus_{i\geq0}P_R)$ has
$\mathbb D_\hbar(M)\simeq W$, which is not supported.
If $E$ is perfect over $R$, then
\[
 \mathbb D_\hbar(I(E))\simeq I(E^\vee),
 \qquad E^\vee=\operatorname{RHom}_R(E,R).
\]
\end{proposition}
\begin{proof}
Dualizing~\eqref{mht:eq:normal-free-resolution} gives
$[\widehat D_R\xrightarrow{R_u}\widehat D_R]$ in degrees
$0,1$. Here $R_u$ is right multiplication by $u$, a
left-module map. In the PBW order $\partial^ju^i$, it
is an injective shift of $i$, with cokernel
$\widehat P_R^t$. Thus the normalized dual of $P_R$
is $\widehat P_R$.

The complete target allows removal of source completion
inside this derived Hom. Mapping out of the direct sum
gives $W$. The source $M$ is flat, and its reductions
$\bigoplus_iP_{R_r}$ are supported. The preceding
proposition proves the failure for its dual.
For perfect $E$, finite projective evaluation commutes
with this calculation and introduces no infinite product.
It gives the second formula.
\end{proof}

Keeping the finite-sum duals as a pro diagram preserves
their support. It cannot be replaced by an ordinary
supported realization preserving every perfect test.
Indeed, already over $k=\mathbb C$, such a realization
$Z$ would have natural equivalences
\[
 \operatorname{RHom}_{D_k}(F,Z)\simeq
 \mathbf R\varprojlim_n
             \operatorname{RHom}_{D_k}(F,P_k^{\oplus n})
 \qquad(F\text{ perfect}).
\]

The free test $F=D_k$ identifies its underlying complex
with $\prod_iP_k$, concentrated in degree zero.
Naturality under the endomorphisms $L_d$ of $D_k$
recovers the right $D_k$-action. The diagonal family
then contradicts support. Restricting to supported tests
would change this requirement: over $R$ they do not detect
$L$, and over $k$ they do not detect its reduction.

Perfect scalar grounds also need a relative finiteness
condition if relative tensors are allowed. For
$B=R[\epsilon]/(\epsilon^2)$ and the character module
$R=B/(\epsilon)$, the free resolution has differential
$\epsilon$ at every step. Tensoring with $R$ makes all
differentials zero. Consequently
$H^{-n}(\Lambda(R\otimes_B^{\mathbb L}R))=R$ for every
$n\geq0$, although each input is scalar-perfect.
The output remains supported after induction but is no
longer scalar-perfect.

These calculations specify an algebraic normal model with
reconstruction and additive convolution, together with
its limits of closure. Identifying this convolution with
boundary fusion requires the geometric correspondence,
tangential factors, unit and supported collision maps.
Its unit $\widehat P_R$, the scalar unit $R$, the operator
unit $1\in D_R$ and a prepared physical state have distinct
roles. The physical comparison must also decide which
probes are allowed: retaining a pro diagram and taking
its supported limit preserve different test functors.
